% =============================================================================
% PRESENTACIÓN REDVAR - Proceso de bioconversión con Hermetia illucens
% Compilación: xelatex presentacion-redvar.tex
% =============================================================================
\documentclass[aspectratio=169, 11pt]{beamer}

% --- Fuentes (XeLaTeX) ---
\usepackage{fontspec}
\setmainfont{TeX Gyre Termes}
\setsansfont{TeX Gyre Heros}
\usepackage[spanish,es-nodecimaldot]{babel}
\usepackage{amsmath, amssymb}
\usepackage{booktabs}
\usepackage{graphicx}

% --- Tema y colores ---
\usetheme{Madrid}
\usecolortheme{beaver}
\definecolor{unalRojo}{RGB}{148, 36, 36}
\definecolor{verdeOK}{RGB}{34, 139, 34}
\setbeamercolor{structure}{fg=unalRojo}
\setbeamercolor{block title}{bg=unalRojo!15, fg=unalRojo}
\setbeamercolor{block body}{bg=unalRojo!5}

% --- Metadatos ---
\title[Bioconversión BSF]{Del residuo al recurso: bioconversión con \textit{Hermetia illucens}}
\subtitle{Proceso, monitoreo ambiental y revalorización}
\author{John Jairo Leal Gómez}
\institute[UNAL]{Doctorado en Estudios Ambientales\\Universidad Nacional de Colombia -- Sede Palmira}
\date{REDVAR -- Evento de Revalorización de Residuos\\Noviembre 2026}
\titlegraphic{\includegraphics[width=1.5cm]{../tesis/00Figuras/00f00EscudoUN2016.jpg}}

\newcommand{\BSF}{\textit{Hermetia illucens}}
\setbeamertemplate{navigation symbols}{}
\setbeamertemplate{footline}[frame number]

\begin{document}

% =============================================================================
% APERTURA - EL PROBLEMA
% =============================================================================

% --- 1. Portada ---
\begin{frame}[plain]
    \titlepage
\end{frame}

% --- 2. El problema ---
\begin{frame}{Residuos orgánicos: fuente de GEI y oportunidad perdida}
% TODO: cifras GEI Colombia/Valle, descomposición anaerobia en vertederos
\end{frame}

% =============================================================================
% EL PROCESO
% =============================================================================

% --- 3. La apuesta ---
\begin{frame}{La apuesta: bioconversión con \BSF{}}
% TODO: qué es, velocidad (días), productos, por qué BSF
\end{frame}

% --- 4. El problema operativo ---
\begin{frame}{El proceso opera como caja negra}
% TODO: sin medición no hay gestión; anoxia local, pérdida de biomasa
\end{frame}

% --- 5. Medir para saber ---
\begin{frame}{Medir para saber: nodo IoT de bajo costo}
% TODO: ESP32 + NDIR (CO2) + electroquímico (CH4), T/HR, calibración
\end{frame}

% --- 6. Modelar para entender ---
\begin{frame}{Modelar para entender: curva esperada vs. medición}
% TODO: modelo DEB (Eriksen 2022), CO2 esperado como referencia
\end{frame}

% --- 7. Detectar para actuar ---
\begin{frame}{Detectar para actuar: alerta de anaerobiosis}
% TODO: residuo operativo Y_sensor - Y_modelo, correlación con CH4
\end{frame}

% --- 8. Lo que muestran los datos ---
\begin{frame}{Lo que muestran los ciclos reales}
% TODO: gráficas de datos experimentales, picos CH4, efecto dieta
\end{frame}

% =============================================================================
% REVALORIZACIÓN
% =============================================================================

% --- 9. Doble salida ---
\begin{frame}{La doble salida: biomasa + frass}
% TODO: biomasa 37-56% proteína, frass N-P-K
\end{frame}

% --- 10. Frass ---
\begin{frame}{Frass: de excreta a fertilizante orgánico}
% TODO: composición, potencial agrícola regional
\end{frame}

% --- 11. Proteína ---
\begin{frame}{Proteína animal: alternativa a harina de pescado}
% TODO: acuicultura y avicultura
\end{frame}

% --- 12. Cadena de valor ---
\begin{frame}{Residuos agroindustriales del Valle: cadena de valor}
% TODO: pulpa de café, bagazo, orujo; del residuo al mercado
\end{frame}

% =============================================================================
% CIERRE
% =============================================================================

% --- 13. Cierre ---
\begin{frame}{Medir, valorizar, escalar}
% TODO: mensaje final + contacto
\end{frame}

\end{document}
